\documentclass[10pt,aspectratio=169]{beamer}
\usetheme[progressbar=frametitle, numbering=fraction, block=fill]{metropolis}
\usepackage[T1]{fontenc}
\usepackage[utf8]{inputenc}
\usepackage[brazilian]{babel}
\usepackage{booktabs,tabularx,array}
\usepackage{amsmath,amssymb}
\usepackage{pgfplots,pgfplotstable}
\usepgfplotslibrary{statistics}
\usepackage{tabuleiro}
% graficos.tex -- figuras pgfplots compartilhadas por relatorio.tex e slides.tex.
% Todas leem os arquivos de dados/ gerados por exportar_dados.py.
% Cada macro aceita um argumento opcional com opções extras de eixo (largura, altura...).

\pgfplotsset{
  compat=1.18,
  tde/.style={
    tick label style={font=\scriptsize},
    label style={font=\footnotesize},
    legend style={font=\scriptsize, draw=none, fill=white, fill opacity=0.85, text opacity=1},
    legend cell align=left,
    axis line style={black!60},
    major tick style={black!60},
    grid style={black!10},
  },
  adv ale/.style={color=corAle, mark=none, thick},
  adv d2/.style={color=corDdois, mark=none, thick},
  adv d4/.style={color=corDquatro, mark=none, thick},
  adv d6/.style={color=corDseis, mark=none, thick},
}
\pgfplotstableread{dados/resultados.dat}\tabResultados
\pgfplotstableread{dados/porlado.dat}\tabPorLado
\pgfplotstableread{dados/caixas.dat}\tabCaixas
\pgfplotstableread{dados/hist_nos.dat}\tabHistNos
\pgfplotstableread{dados/ordenacao.dat}\tabOrdenacao
\pgfplotstableread{dados/duracao.dat}\tabDuracao
\pgfplotstableread{dados/zugzwang.dat}\tabZugzwang

\def\rotulosAdv{Aleatório,AlfaBeta-d2,AlfaBeta-d4,AlfaBeta-d6}

% ---------- V/E/D empilhados + taxa de vitória com IC 95% de Wilson ----------
\newcommand{\graficoResultados}[1][]{%
\begin{tikzpicture}
  \begin{axis}[tde, xbar stacked, bar width=9pt, width=0.9\linewidth, height=4.6cm,
      xmin=0, xmax=1, ymin=-0.6, ymax=3.6, y dir=reverse,
      ytick={0,1,2,3}, yticklabels/.expand once=\rotulosAdv,
      xtick={0,0.25,0.5,0.75,1}, xticklabels={0\,\%,25\,\%,50\,\%,75\,\%,100\,\%},
      xmajorgrids, xlabel={fração das partidas do AlfaBeta-d8},
      legend style={at={(0.5,1.02)}, anchor=south, legend columns=4},
      area legend, #1]
    \addplot[fill=corV, draw=none] table[x=vfrac, y=i] \tabResultados;
    \addplot[fill=corE, draw=none] table[x=efrac, y=i] \tabResultados;
    \addplot[fill=corD!80, draw=none] table[x=dfrac, y=i] \tabResultados;
    \addlegendentry{vitória}\addlegendentry{empate}\addlegendentry{derrota}
    \addlegendimage{black, mark=|, only marks, line width=1pt}\addlegendentry{IC 95\,\% da vitória}
  \end{axis}
  \begin{axis}[width=0.9\linewidth, height=4.6cm, xmin=0, xmax=1, ymin=-0.6, ymax=3.6, y dir=reverse,
      hide axis, #1]
    \addplot[black, line width=1pt, mark=none, only marks,
      error bars/.cd, x dir=both, x explicit, error bar style={line width=1pt},
      error mark options={rotate=90, mark size=3pt, line width=1pt}]
      table[x=taxa, y=i, x error plus expr=\thisrow{hi}-\thisrow{taxa},
            x error minus expr=\thisrow{taxa}-\thisrow{lo}] \tabResultados;
  \end{axis}
\end{tikzpicture}}

% ---------- distribuição de nós por lance (caixas em escala log) ----------
% caixa = quartis; linha = mediana; bigodes = mínimo e p95; × = máximo; ◆ = média
\newcommand{\graficoCaixas}[1][]{%
\begin{tikzpicture}
  \begin{axis}[tde, ymode=log, width=0.9\linewidth, height=6cm, boxplot/draw direction=y,
      xtick={1,2,3,4}, xticklabels/.expand once=\rotulosAdv, xmin=0.4, xmax=4.6,
      ymin=1, ymax=4e6, ymajorgrids, log basis y=10, ytick={1,10,100,1000,10000,100000,1000000},
      ylabel={nós visitados por lance do d8}, #1]
    \pgfplotsinvokeforeach{0,1,2,3}{
      \pgfplotstablegetelem{##1}{min}\of\tabCaixas\let\cxA\pgfplotsretval
      \pgfplotstablegetelem{##1}{q1}\of\tabCaixas\let\cxB\pgfplotsretval
      \pgfplotstablegetelem{##1}{med}\of\tabCaixas\let\cxC\pgfplotsretval
      \pgfplotstablegetelem{##1}{q3}\of\tabCaixas\let\cxD\pgfplotsretval
      \pgfplotstablegetelem{##1}{p95}\of\tabCaixas\let\cxE\pgfplotsretval
      \edef\cxPlot{\noexpand\addplot[fill=corTab!15, draw=corTab, boxplot prepared={draw position=##1+1,
        lower whisker=\cxA, lower quartile=\cxB, median=\cxC, upper quartile=\cxD, upper whisker=\cxE}]
        coordinates {};}\cxPlot
    }
    \addplot[only marks, mark=x, mark size=3pt, corD, thick] table[x expr=\thisrow{i}+1, y=max] \tabCaixas;
    \addplot[only marks, mark=diamond*, mark size=2.5pt, black] table[x expr=\thisrow{i}+1, y=media] \tabCaixas;
  \end{axis}
\end{tikzpicture}}

% ---------- histograma de nós por lance (bins logarítmicos) ----------
\newcommand{\graficoHistNos}[1][]{%
\begin{tikzpicture}
  \begin{axis}[tde, xmode=log, width=0.9\linewidth, height=5cm, xmin=1, xmax=3.2e6,
      ymin=0, xlabel={nós visitados por lance do d8 (4 faixas por década)},
      ylabel={fração dos lances}, ymajorgrids,
      yticklabel={\pgfmathparse{100*\tick}\pgfmathprintnumber[fixed,precision=0]{\pgfmathresult}\,\%},
      legend pos=north west, #1]
    \addplot[adv ale, const plot] table[x=x, y=ale] \tabHistNos; \addlegendentry{vs Aleatório}
    \addplot[adv d2, const plot] table[x=x, y=d2] \tabHistNos; \addlegendentry{vs d2}
    \addplot[adv d4, const plot] table[x=x, y=d4] \tabHistNos; \addlegendentry{vs d4}
    \addplot[adv d6, const plot] table[x=x, y=d6] \tabHistNos; \addlegendentry{vs d6}
  \end{axis}
\end{tikzpicture}}

% ---------- ordenação: com vs sem, mesma posição ----------
\newcommand{\graficoOrdenacao}[1][]{%
\begin{tikzpicture}
  \begin{axis}[tde, xmode=log, ymode=log, width=0.62\linewidth, height=0.62\linewidth,
      xmin=10, xmax=1e6, ymin=10, ymax=1e6, grid=major,
      xlabel={nós sem ordenação}, ylabel={nós com centro primeiro},
      legend pos=north west, #1]
    \addplot[black!50, dashed, domain=10:1e6, samples=2, forget plot] {x};
    \addplot[scatter, only marks, scatter src=explicit symbolic,
      scatter/classes={melhor={mark=*, mark size=1.2pt, corTab, fill opacity=0.6, draw opacity=0.6},
                       pior={mark=*, mark size=1.6pt, corD}}]
      table[x=sem, y=com, meta=classe] \tabOrdenacao;
    \legend{ordenação ajudou, ordenação custou mais}
    \node[font=\scriptsize, black!60, rotate=45, anchor=south] at (axis cs:2e5,2.6e5) {$y = x$};
  \end{axis}
\end{tikzpicture}}

% ---------- duração das partidas ----------
\newcommand{\graficoDuracao}[1][]{%
\begin{tikzpicture}
  \begin{axis}[tde, width=0.9\linewidth, height=4.6cm, xmin=5, xmax=45, ymin=0,
      xtick={5,10,...,45}, xlabel={peças no tabuleiro no fim da partida (faixas de 5)},
      ylabel={fração das partidas}, ymajorgrids,
      yticklabel={\pgfmathparse{100*\tick}\pgfmathprintnumber[fixed,precision=0]{\pgfmathresult}\,\%},
      legend style={at={(0.5,1.03)}, anchor=south, legend columns=4}, #1]
    \addplot[adv ale, const plot] table[x=x, y=ale] \tabDuracao; \addlegendentry{vs Aleatório}
    \addplot[adv d2, const plot] table[x=x, y=d2] \tabDuracao; \addlegendentry{vs d2}
    \addplot[adv d4, const plot] table[x=x, y=d4] \tabDuracao; \addlegendentry{vs d4}
    \addplot[adv d6, const plot] table[x=x, y=d6] \tabDuracao; \addlegendentry{vs d6}
  \end{axis}
\end{tikzpicture}}

% ---------- exemplo de zugzwang: AVAL e valor da busca d8 a cada lance do d8 ----------
\newcommand{\graficoZugzwang}[1][]{%
\begin{tikzpicture}
  \begin{axis}[tde, width=0.9\linewidth, height=5cm, xmin=2, xmax=35, ymin=-16, ymax=42,
      xlabel={peças no tabuleiro quando o d8 (O) decide}, ylabel={ponto de vista de X (MAX)},
      ymajorgrids, legend style={at={(0.02,0.98)}, anchor=north west},
      ytick={-15,-10,-5,0,5,10,15,30}, yticklabels={$-15$,$-10$,$-5$,0,5,10,15,{+965}}, #1]
    \fill[corD!8] (axis cs:\dado{zz-perdido-pecas}-1,-16) rectangle (axis cs:35,42);
    \draw[black!40] (axis cs:2,0) -- (axis cs:35,0);
    \addplot[corTab, thick, mark=*, mark size=1.5pt] table[x=pecas, y=valor] \tabZugzwang;
    \addlegendentry{valor da busca d8}
    \addplot[corO!80!black, thick, dashed, mark=o, mark size=1.5pt] table[x=pecas, y=aval] \tabZugzwang;
    \addlegendentry{AVAL da posição}
    \node[font=\scriptsize, anchor=north, text=corD, align=center] at (axis cs:31,25) {derrota\\já forçada};
  \end{axis}
\end{tikzpicture}}


\setbeamercolor{progress bar}{fg=corTab}
\setbeamercolor{alerted text}{fg=corD}
\setbeamercolor{frametitle}{bg=corTabEsc}
\newcommand{\fn}[1]{\text{\footnotesize\MakeUppercase{#1}}}
\newcommand{\arq}[1]{\texttt{#1}}
\newcommand{\tempo}[1]{\hfill{\scriptsize\color{black!45}#1}}
\newcommand{\captura}[2]{%
  \IfFileExists{figs/#1.png}{\includegraphics[width=#2]{figs/#1.png}}{%
    \fbox{\parbox[c][3cm][c]{\dimexpr#2-2\fboxsep-2\fboxrule\relax}{\centering\small
      \arq{figs/#1.png} ausente: \arq{make capturas}}}}}

\title{TDE2 --- Busca Competitiva}
\subtitle{Um agente alfa-beta para o Lig 4}
\author{\small Guilherme Arcanjo de Morais · Lucas da Costa Paula\\ Renan da Silva Andrade · Renan Americo Herculano}
\institute{\footnotesize PUCPR · Inteligência Artificial (PCO118A) · Prof. Carlos Costa Cedro · RA2}
\date{}

\begin{document}
\maketitle

% =====================================================================
\section{O problema \tempo{1 min}}

\begin{frame}{Lig 4 (Connect Four), 6×7}
\begin{columns}[c]
\column{0.58\textwidth}
\begin{itemize}
  \item X e O soltam peças em colunas; a peça cai até a casa livre mais baixa.
  \item Vence quem alinhar 4. Tabuleiro cheio é empate.
\end{itemize}
\medskip
\small
\begin{tabular}{@{}ll@{}}
\toprule
Dois jogadores & sim \\
Soma zero & sim: $U(s,O) = -U(s,X)$ \\
Informação perfeita & sim \\
Determinístico & sim \\
Finito & sim: no máximo 42 jogadas \\
\bottomrule
\end{tabular}
\column{0.38\textwidth}\centering
\tabuleiro[x=6mm,y=6mm]{\tabZzFinal}[\tabZzFinalVence]\\[4pt]
{\scriptsize uma partida do experimento: X vence na diagonal}
\end{columns}
\medskip
\alert{Por que o método:} pequeno para implementar e grande demais para o Minimax puro.
\end{frame}

% =====================================================================
\section{Modelagem \tempo{3 min}}

\begin{frame}{Estado, ações e transição}
\begin{columns}[t]
\column{0.55\textwidth}
\textbf{Estado} $s = (x, o, \mathit{alturas}, \mathit{vez})$
\begin{itemize}\small
  \item $x, o$: bitboards de 49 bits; casa $(c, r)$ é o bit $7c + r$
  \item 1 bit sentinela (sempre 0) no topo de cada coluna
  \item $\fn{Jogador}(s) = s.\mathit{vez}$; $s_0$: vazio, X (MAX) começa
\end{itemize}
\textbf{\fn{Ações}}$(s) = \{c : \mathit{alturas}[c] < 6\}$, ou $\emptyset$ se terminal\\[3pt]
\textbf{\fn{Resultado}}$(s,c)$: liga o bit $7c + \mathit{alturas}[c]$ de quem joga e passa a vez\\[3pt]
\textbf{\fn{Teste-Terminal}}: 4 em linha ou 42 peças;\\
\quad $b \,\&\, (b{\gg}d) \,\&\, (b{\gg}2d) \,\&\, (b{\gg}3d) \ne 0$, $d \in \{1,7,6,8\}$
\column{0.42\textwidth}
\textbf{\fn{Utilidade}}$(s,\text{MAX})$: um único número
\[
\begin{cases}
+(1000-n) & \text{X venceu}\\
-(1000-n) & \text{O venceu}\\
0 & \text{empate}
\end{cases}
\]
{\small $n$ = peças no tabuleiro: vitória mais rápida vale mais.}
\medskip

\begin{block}{Corte nunca parece terminal}
$|\fn{Aval}| \le 900 < 958 \le |\fn{Utilidade}|$ de qualquer vitória
\end{block}
\end{columns}
\end{frame}

\begin{frame}{Tamanho do espaço e o que a representação proíbe}
\begin{columns}[t]
\column{0.48\textwidth}
\textbf{b e m medidos} (nos estados em que o agente jogou)
\begin{itemize}\small
  \item $b \approx \dado{b-medio}$ (7 no início), $m \approx \dado{m-medio}$ (máx.\ 42)
  \item $7^{42} \approx \dado{bm-max}$ folhas (limite)
  \item $\dado{b-medio}^{\,\dado{m-medio}} \approx \dado{bm-medido}$
  \item $\approx 4{,}5\times10^{12}$ posições legais (Tromp)
\end{itemize}
\alert{A $10^6$ nós/s, $10^{20}$ folhas levariam milhões de anos}: alfa-beta + corte de profundidade.
\column{0.5\textwidth}
\textbf{Impossível por construção}
\begin{itemize}\small
  \item peça flutuando: só entra na casa $(c, \mathit{alturas}[c])$
  \item jogar em coluna cheia: fora de \fn{Ações}
  \item jogar duas vezes: \fn{Resultado} sempre passa a vez
  \item casa dos dois: só a casa livre recebe peça
  \item alinhamento entre colunas: o bit sentinela zera o \texttt{\&}
\end{itemize}
\end{columns}
\end{frame}

\begin{frame}{Implementação: avaliação, ordenação e parada}
\begin{columns}[c]
\column{0.6\textwidth}
\small
$\fn{Aval}(s) = \sum_{\text{janelas}} [5\cdot\text{(3+1 vazia)} + 2\cdot\text{(2+2 vazias)}]_{X - O}$\\
\hspace*{4.2em}$+\; 3\cdot(\text{centro}_X - \text{centro}_O)$, limitada a $\pm 900$
\begin{itemize}
  \item janela com peças dos dois lados vale 0
  \item pesos à mão, iguais para todos os agentes
\end{itemize}
\textbf{Ordenação}: centro primeiro, 3, 2, 4, 1, 5, 0, 6 (custo zero)\\[3pt]
\textbf{Parada}: profundidade fixa $d$. Com relógio, nós e jogadas dependeriam da máquina\\[3pt]
\textbf{Custo}: cada chamada recursiva = 1 nó visitado (raiz, cortes e terminais incluídos)
\column{0.36\textwidth}\centering
\mapaJanelas[x=5mm,y=5mm]\\[2pt]
{\scriptsize janelas de 4 por casa: centro \dado{janelas-max-casa}, canto \dado{janelas-canto}}
\end{columns}
\end{frame}

% =====================================================================
\section{Execução ao vivo \tempo{3 min}}

\begin{frame}[fragile]{Visualizador: \texttt{python3 visualizador.py}}
\begin{columns}[c]
\column{0.68\textwidth}\centering
\captura{jogar}{\linewidth}
\column{0.3\textwidth}\small
\begin{itemize}
  \item \alert{contador de nós ao vivo} enquanto o agente pensa
  \item valor de cada coluna na raiz (exato ou limite $\le/\ge$)
  \item decomposição da \fn{Aval}
  \item reproduz qualquer partida do experimento e confere com \arq{resultados.json}
\end{itemize}
\end{columns}
\end{frame}

\begin{frame}{Árvore alfa-beta passo a passo}
\begin{columns}[c]
\column{0.68\textwidth}\centering
\captura{arvore}{\linewidth}
\column{0.3\textwidth}\small
\begin{itemize}
  \item $\alpha$ e $\beta$ no caminho atual
  \item subárvores podadas em cinza
  \item nós do alfa-beta $\times$ nós do Minimax
  \item mesmo valor na raiz
\end{itemize}
\end{columns}
\end{frame}

\begin{frame}{Corretude: alfa-beta = Minimax em profundidade completa}
\small
\begin{tabularx}{\linewidth}{Xrrc}
\toprule
Caso (\arq{verificacao.py}: \alert{\dado{verificacao-status}}) & Minimax & Alfa-beta & Diverg. \\
\midrule
4×4, $K=3$, a partir de $s_0$ & \dado{t1-mm-valor}, \dado{t1-mm-nos} nós & \dado{t1-ab-valor}, \dado{t1-ab-nos} nós & 0 \\
200 posições 3×4, $K=3$ & \dado{t2-mm} & \dado{t2-ab} (\dado{t2-pct}) & 0 \\
100 finais 6×7 (8 vazias) & \dado{t3-mm} & \dado{t3-ab} (\dado{t3-pct}) & 0 \\
100 posições 6×7, $d=4$, valor e jogada & \dado{t4-mm} & \dado{t4-ab} (\dado{t4-pct}) & 0 \\
\ifnum\dado{teste5}=1
\dado{t5-comparacoes} casos: busca instrumentada do visualizador & \multicolumn{2}{c}{mesmo valor, jogada e nós} & 0 \\
\fi
\bottomrule
\end{tabularx}
\medskip

$991 = 1000 - 9$: no 4×4 com $K=3$, X vence com jogo perfeito na 9.ª peça.\\
Pelo Teorema 1 (Aula 7), a igualdade de valor é a evidência de que o alfa-beta não poda demais.
\end{frame}

% =====================================================================
\section{Comparação \tempo{3 min}}

\begin{frame}{Protocolo}
\begin{itemize}
  \item \textbf{Avaliado}: AlfaBeta-d8. \textbf{Adversários}: Aleatório (linha de base) e o próprio agente com d2, d4, d6
  \item \textbf{Custo}: nós visitados por lance; tempo só informativo
  \item \textbf{Variação}: 49 aberturas de 2 jogadas forçadas $\times$ 2 lados = 98 partidas por confronto;
        aleatório com semente = nº da partida, alternando quem começa
  \item \textbf{Calibração simétrica}: mesmo código, \fn{Aval}, pesos e ordenação; só $d$ muda
  \item \dado{total-partidas} partidas, determinísticas: \arq{python3 experimento.py} reproduz tudo
\end{itemize}
\end{frame}

\begin{frame}{Resultados do AlfaBeta-d8}
\begin{columns}[c]
\column{0.5\textwidth}
\footnotesize\setlength{\tabcolsep}{4pt}
\begin{tabular}{@{}lccr@{}}
\toprule
Adversário & V/E/D & Vitória (IC 95\,\%) & Pontos \\
\midrule
Aleatório & \dado{ale-V}/\dado{ale-E}/\dado{ale-D} & \dado{ale-taxa} (\dado{ale-lo} a \dado{ale-hi}) & \dado{ale-pontos} \\
d2 & \dado{d2-V}/\dado{d2-E}/\dado{d2-D} & \dado{d2-taxa} (\dado{d2-lo} a \dado{d2-hi}) & \dado{d2-pontos} \\
d4 & \dado{d4-V}/\dado{d4-E}/\dado{d4-D} & \dado{d4-taxa} (\dado{d4-lo} a \dado{d4-hi}) & \dado{d4-pontos} \\
d6 & \dado{d6-V}/\dado{d6-E}/\dado{d6-D} & \dado{d6-taxa} (\dado{d6-lo} a \dado{d6-hi}) & \dado{d6-pontos} \\
\bottomrule
\end{tabular}
\medskip

\small 98 partidas por adversário alfa-beta, 100 contra o aleatório.\\
Contra o d6 o IC fica abaixo de 50\,\%: \alert{não dá para afirmar d8 > d6}.
\column{0.47\textwidth}
\graficoResultados[width=0.9\linewidth, height=5cm, yticklabels={Aleat.,d2,d4,d6},
  legend style={legend columns=2}]
\end{columns}
\end{frame}

\begin{frame}{Custo por lance: a distribuição inteira}
\begin{columns}[c]
\column{0.52\textwidth}
\scriptsize\setlength{\tabcolsep}{4pt}
\begin{tabular}{@{}lrrrrr@{}}
\toprule
Nós/lance & Média & Mediana & p95 & Máx. & Adv. \\
\midrule
vs Aleat. & \dado{ale-media} & \dado{ale-mediana} & \dado{ale-p95} & \dado{ale-max} & \dado{ale-advmedia} \\
vs d2 & \dado{d2-media} & \dado{d2-mediana} & \dado{d2-p95} & \dado{d2-max} & \dado{d2-advmedia} \\
vs d4 & \dado{d4-media} & \dado{d4-mediana} & \dado{d4-p95} & \dado{d4-max} & \dado{d4-advmedia} \\
vs d6 & \dado{d6-media} & \dado{d6-mediana} & \dado{d6-p95} & \dado{d6-max} & \dado{d6-advmedia} \\
\bottomrule
\end{tabular}
\medskip

\small Cauda pesada: a média fica 2 a 3 vezes acima da mediana, e o máximo passa de 1 milhão de nós.
\column{0.46\textwidth}
\graficoCaixas[width=\linewidth, height=5.4cm, xticklabels={Aleat.,d2,d4,d6}, ylabel={nós por lance}]
\end{columns}
\end{frame}

\begin{frame}{Ordenação centro-primeiro}
\begin{columns}[c]
\column{0.45\textwidth}\centering
\graficoOrdenacao[width=0.95\linewidth, height=0.95\linewidth]
\column{0.53\textwidth}
\begin{itemize}\small
  \item mesmas \dado{ord-posicoes} posições, d8
  \item média \dado{ord-com-media} com $\times$ \dado{ord-sem-media} sem (\(-\)\dado{ord-reducao-media})
  \item mediana \dado{ord-com-mediana} $\times$ \dado{ord-sem-mediana}
  \item valor da raiz idêntico: \dado{ord-mesmos}
  \item \alert{em \dado{ord-pior} posições custou mais}, até \dado{ord-razao-max}$\times$
\end{itemize}
\end{columns}
\end{frame}

% =====================================================================
\section{O que não funcionou \tempo{2 min}}

\begin{frame}{A \fn{Aval} ignora a paridade das ameaças}
\vspace*{2mm}
\begin{columns}[c]
\column{0.5\textwidth}\centering
\tabuleiro[x=3.2mm,y=3.2mm]{\tabZzAntes}[][\colZzAntes]\hspace{2mm}%
\tabuleiro[x=3.2mm,y=3.2mm]{\tabZzPerdido}[][\colZzPerdido]\\[2pt]
{\scriptsize \dado{zz-antes-pecas} peças: valor \dado{zz-antes-valor} \hspace{6mm} \dado{zz-perdido-pecas} peças: \dado{zz-perdido-valor}}
\medskip

\graficoZugzwang[width=\linewidth, height=3.7cm, xlabel={peças quando o d8 (O) decide}, ylabel={}]
\column{0.48\textwidth}
\small
\begin{itemize}
  \item o d8 perdeu \dado{d2-D} partidas até para o d2, \dado{d2-perdidas-O} delas como O, todas com
        \dado{d2-perdidas-O-min} a \dado{d2-perdidas-O-max} peças
  \item partida nº \dado{zz-id}: com \dado{zz-antes-pecas} peças a busca ainda via \dado{zz-antes-valor};
        com \dado{zz-perdido-pecas}, já via a derrota forçada
  \item ameaças em linha ímpar ou par decidem o final, e a \fn{Aval} conta as duas igual
\end{itemize}
\end{columns}
\end{frame}

\begin{frame}{Garantias que o sistema não tem}
\begin{itemize}
  \item \textbf{Não é ótimo}: só $|v| \ge 958$ é resultado provado; o resto é estimativa
  \item \textbf{Efeito horizonte}: não há busca de quiescência
  \item \textbf{Custo sem teto}: profundidade fixa; pior lance medido \dado{d2-max} nós
  \item \textbf{Pesos não ajustados}: (5, 2, 3) escolhidos à mão
  \item \textbf{Validade restrita}: contra alfa-beta, só as 49 aberturas testadas
  \item \textbf{Determinístico}: um adversário pode repetir uma linha vencedora
  \item \textbf{Retornos decrescentes}: d8 $\times$ d6 indistinguível com 98 partidas
\end{itemize}
\end{frame}

\end{document}
